% =====================================================================
%  Chapter 15 opener -- "Chapter at a Glance", three reading paths + the
%  ladder staircase, THE running example (the same 10-node tree T as
%  Ch.12/Ch.13/Ch.14, edges 1-2,1-3,2-4,2-5,3-6,3-7,5-8,7-9,9-10) re-drawn
%  with a rerooting (down, up) pair overlay, the "who owns what" comparison
%  (Ch.13 fixed-tuple cookbook vs Ch.14 budgeted/game DPs vs THIS chapter's
%  rerooting-as-a-theory vs Ch.16 LCA vs Ch.17 DSU vs Ch.24 centroid),
%  roadmap, recognition guide, constraint->tool table, notation.
%  Exports fig:c20-runtree-pair and fig:c20-taxonomy.
%  All overlays below are hand-computed on T; every later section re-uses
%  exactly these numbers.  For the anchor fold of 15.1 (CF 1187E), down is
%  the Ch.13 subtree-size fold and up is the parent-side piece the edge-
%  transfer rule manufactures; the pair sums to the every-root answer F(v).
% =====================================================================

\begin{chaptertheorybox}{Chapter at a Glance}
Chapter~13 rooted the tree \emph{once}, folded a constant tuple up from the
children, and read the answer at that single root. Chapter~14 kept the root
fixed but gave the tuple an index. This chapter finally drops the fixed-root
assumption: the answer must be reported \emph{for every choice of root}, and
we may not afford to re-run the DFS $n$ times.  The whole subject collapses
to one observation.

Moving the root across a single edge is a \term{constant-size update}.  Split
every node's answer into a \emph{down} part (what its own subtree contributes,
Ch.~13's fold) and an \emph{up} part (what the rest of the tree, reached
through its parent, contributes).  The down pass is exactly
\thmref{c10-subtree-template}, reused verbatim.  The up pass is the new
material: a second DFS, top-down, that transfers the answer from a parent
$v$ to each child $c$ by the \term{edge-transfer rule}
$\mathit{up}[c] = \Phi\!\operatorname{-shift}\!\bigl(\text{combine}(v\text{'s self},
\mathit{up}[v], \text{all of }c\text{'s siblings})\bigr)$
(\thmref{c20-edge-transfer}).  One DFS down, one DFS up, and every node holds
its own rooted answer.

\textbf{You will learn to}
\begin{itemize}[nosep]
  \item state the two-DFS template -- a \emph{down} fold and an \emph{up}
        fold -- and join them into $\mathit{dp}[v]=\Phi(\mathit{down}[v],
        \mathit{up}[v])$ for every root in $\bigO{n}$; \hfill\lv{2}--\lv{3}
  \item derive the edge-transfer rule for a concrete fold and prove it,
        then read CF~1187E's $\mathit{up}[c]=\mathit{up}[v]+n-2\size(c)$
        as the special case; \hfill\lv{3}
  \item spot a \emph{non-invertible} combine ($\max$, signed product,
        $\min$, $\mathit{or}$) and replace "subtract the child" with a
        prefix/suffix merge over $v$'s children list -- still $\bigO{n}$
        total; \hfill\lv{4}
  \item carry a pair-valued or count-valued state through both passes and
        still land $\bigO{n}$ (CF~1156D, CF~960E); \hfill\lv{4}--\lv{5}
  \item reroot a game/DP \emph{verdict} for every root, tying Ch.~14's
        Grundy theory to the two-DFS template (CF~1626E); \hfill\lv{5}
  \item combine everything on the Summit CF~1498F: reroot a Grundy value
        across every node, discharging Ch.~14's "rerooting the cutting game
        over every ground" promise. \hfill\lv{5}--\lv{6}
\end{itemize}

\textbf{Prerequisites.} Chapter~13 (\chref{c10}) for the constant-tuple
child-fold \thmref{c10-subtree-template} -- that fold \emph{is} this chapter's
down pass; Chapter~12 (\chref{c09}) for the running tree $T$, $\size(v)$, and
the measuring DFS \coderef{c09-dfs3}; Chapter~14 (\chref{c27}) for the
budgeted merge \thmref{c27-amortize} and the Grundy/Hackenbush theory
\thmref{c27-hackenbush} the Summit reroots. Chapter~4 (\chref{c04}) for the
rooting-by-parent idiom. Nothing from Chapters~16/17/24 (LCA, DSU-on-tree,
centroid) is assumed -- this chapter is the theory they are later contrasted
against.
\end{chaptertheorybox}

% ---------------- Reading paths + ladder position ----------------
\begin{chaptertheorybox}{How to Read This Chapter (three paths)}
Level badges $\lv1$--$\lv6$: $\lv2$ is applying the two-DFS template to a
fold whose edge-transfer you are handed; $\lv3$ is deriving the edge-transfer
rule yourself and seeding the up pass correctly; $\lv4$ needs \emph{two} ideas
composed (a pair- or count-valued state, or a non-invertible combine fixed by
a prefix/suffix merge); $\lv5$ is a rerooted game verdict where the state is a
$\operatorname{mex}$; $\lv6$ is reading the two-DFS template as a reusable
black box and knowing when centroid/DSU take over.

\textbf{Canonical two-DFS-first path ($\lv2$--$\lv{3}$, "I want the template").}
Read \secref{c20-twodfs} end to end: the down/up split, \thmref{c20-edge-transfer},
the up-pass seed $\mathit{up}[\text{root}]=\Phi\text{'s identity}$, then the
CF~1187E solved problem and its trace.  You leave owning the single rule that
generates every problem in the chapter.

\textbf{Prefix-suffix-first path ($\lv4$, "my combine has no inverse").}
Skim \secref{c20-twodfs} for the notation, then go straight to
\secref{c20-prefixsuffix}: why $\max$/signed-product/$\min$ DPs cannot use
"subtract the child", and how a prefix/suffix precompute over $v$'s children
buys the "everything except $c_i$" aggregate in $\bigO{1}$ per child.  CF~960E
is the case study.

\textbf{Games-reroot-first path ($\lv5$--$\lv6$, contest track).}  Read
\secref{c20-updown} (pair/count states) and \secref{c20-counting} (rerooting
a win/lose verdict), solve CF~1156D and CF~1626E, then the mixed set, then the
Summit CF~1498F.  At that point the forward ropes are visible: LCA \emph{queries}
are Chapter~16 (\chref{c11}), DSU-on-tree is Chapter~17 (\chref{c28}), centroid
decomposition is Chapter~24 (\chref{c18}).

The backward ropes: Chapter~13 (\chref{c10}) for the fold that is the down
pass, Chapter~12 (\chref{c09}) for $T$ and $\size$, Chapter~14 (\chref{c27})
for the game theory the Summit reroots.
\end{chaptertheorybox}

\begin{figure}[H]
\centering
\begin{tikzpicture}[diagram, x=9mm, y=7mm]
  \foreach \st/\lab in {1/lv1,2/lv2,3/lv3,4/lv4,5/lv5,6/lv6} {
    \fill[\lab!22, draw=\lab, line width=0.8pt] (\st-1,0) rectangle (\st,\st*0.52);
    \node[font=\scriptsize\bfseries, text=\lab] at (\st-0.5,\st*0.52-0.42) {L\st};}
  \node[annot, anchor=south west, font=\scriptsize]
    at (0.05,0.62) {15.1 down/up split + edge-transfer rule \thmref{c20-edge-transfer}};
  \node[annot, anchor=south west, font=\scriptsize]
    at (1.05,1.14) {15.1 CF 1187E \$F(c)=F(v)+n-2\$sz(c)\$};
  \node[annot, anchor=south west, font=\scriptsize]
    at (2.05,1.66) {15.2 pair/count state through both passes (CF 1156D)};
  \node[annot, anchor=south west, font=\scriptsize, text=lv5]
    at (3.05,2.18) {15.3 non-invertible combine, prefix/suffix (CF 960E)};
  \node[annot, anchor=south west, font=\scriptsize, text=lv5]
    at (4.05,2.70) {15.4 reroot a game verdict (CF 1626E)};
  \node[annot, anchor=south west, font=\scriptsize, text=lv6]
    at (5.05,3.22) {Summit CF 1498F reroot Grundy across every node};
\end{tikzpicture}
\caption*{Chapter~15 on the six-step ladder: the down/up split and the
edge-transfer rule sit at $\lv2$--$\lv3$; the pair/count state of CF~1156D is
$\lv3$--$\lv4$; the non-invertible prefix/suffix merge (CF~960E) and the
rerooted game verdict (CF~1626E) climb to $\lv4$--$\lv5$; the Summit CF~1498F
(a rerooted Sprague--Grundy value) is step~5 read from step~6.  Template path
$=$ steps 2--3; full contest path runs to 6.}
\end{figure}

% ---------------- The running example ----------------
\subsection*{The Running Example: One Ten-Node Tree, an Answer at Every Root}

\emph{One} tree does all the work -- the same $T$ Chapters~12, 13 and~14 used,
now read from \emph{every} root.  Nodes $1,\dots,10$; for the reference
orientation we root at $1$, edges
\[
  1\!-\!2,\; 1\!-\!3,\; 2\!-\!4,\; 2\!-\!5,\; 3\!-\!6,\; 3\!-\!7,\; 5\!-\!8,\; 7\!-\!9,\; 9\!-\!10 .
\]
Child sets in that view: $1\!\to\!\{2,3\}$, $2\!\to\!\{4,5\}$,
$3\!\to\!\{6,7\}$, $5\!\to\!\{8\}$, $7\!\to\!\{9\}$, $9\!\to\!\{10\}$; leaves
$\{4,6,8,10\}$.  Subtree sizes $\size(v)$ are $10,4,5,1,2,1,3,1,2,1$ for
$v=1..10$ (Ch.~12/13's table, reused verbatim) -- and their sum,
$30$, is precisely CF~1187E's answer when the root is $1$.

The overlay in \figref{c20-runtree-pair} is the anchor fold of \secref{c20-twodfs}
(15.1): \term{CF~1187E Tree Painting}, where $F(r)$ is the sum of subtree
sizes when the tree is rooted at $r$, and we want $\max_r F(r)$.  Below each
node we print the pair $(\mathit{down}[v],\mathit{up}[v])$ with
$\mathit{down}[v]=\size(v)$ (the Ch.~13 fold) and $\mathit{up}[v]$ the
parent-side piece the edge-transfer rule manufactures; the answer is their
sum $F(v)$.  Root $1$ has $\mathit{up}[1]=0$ (no nodes above it), and the
rule $\mathit{up}[c]=\mathit{up}[v]+(n-\size(c))$ propagates the rest.

\begin{figure}[H]
\centering
\begin{tikzpicture}[diagram, scale=0.90, transform shape,
    tnode/.style={circle, draw=primaryblue, line width=0.9pt, fill=primaryblue!8,
                  minimum size=8mm, font=\small\bfseries, text=primaryblue!60!black}]
  \node[tnode] (n1) at (0,0)     {$1$};
  \node[tnode] (n2) at (-3,-1.6) {$2$};
  \node[tnode] (n3) at ( 3,-1.6) {$3$};
  \node[tnode] (n4) at (-5.2,-3.2) {$4$};
  \node[tnode] (n5) at (-1.6,-3.2) {$5$};
  \node[tnode] (n6) at ( 1.6,-3.2) {$6$};
  \node[tnode] (n7) at ( 5.2,-3.2) {$7$};
  \node[tnode] (n8) at (-1.6,-4.8) {$8$};
  \node[tnode] (n9) at ( 5.2,-4.8) {$9$};
  \node[tnode] (n10) at ( 5.2,-6.4) {$10$};
  \draw[tedge] (n1)--(n2) (n1)--(n3) (n2)--(n4) (n2)--(n5)
               (n3)--(n6) (n3)--(n7) (n5)--(n8) (n7)--(n9) (n9)--(n10);
  % (down, up) = (size, F-size) badges; answer F below
  \foreach \nd/\dn/\up in {
      n1/10/20, n2/4/28, n3/5/25, n4/1/39, n5/2/36,
      n6/1/37, n7/3/31, n8/1/45, n9/2/38, n10/1/47} {
    \node[font=\tiny\ttfamily, text=warningred!75!black, anchor=north]
      at (\nd.south) {(\dn,\up)};}
  \node[annot, font=\scriptsize, anchor=south west, text=successgreen!50!black]
    at (-0.4,0.55) {(down, up) pair; answer F = down + up};
\end{tikzpicture}
\caption{The running tree $T$ with the CF~1187E \emph{rerooting} overlay.
Below each node is the pair $(\mathit{down}[v],\mathit{up}[v])$; the answer is
their sum $F(v)$.  Sanity: at the reference root $\mathit{down}[1]=\size(1)=10$,
$\mathit{up}[1]=0$, so $F(1)=10+\dots$ reads the full $F(1)=30$ once summed over
the fold (the pair is the per-node decomposition, not a running total).  The
maximum is $F(10)=1+47=48$: rooting at the leaf $10$ maximises the sum of
subtree sizes.  Re-derive every $\mathit{up}$ value in \secref{c20-twodfs} from
$\mathit{up}[c]=\mathit{up}[v]+(n-\size(c))$; the whole table is reproduced as
\exref{c20-twodfs-run}.}
\label{fig:c20-runtree-pair}
\end{figure}

Four angles, four hand-computed overlays -- each lives in the section named on
its row and never changes:

\begin{center}
\small
\renewcommand{\arraystretch}{1.25}
\begin{tabular}{@{}p{0.26\linewidth} p{0.46\linewidth} l@{}}
\toprule
\textbf{Angle on $T$} & \textbf{What comes out} & \textbf{Traced in} \\
\midrule
edge-transfer on $\size$ & $F(v)=[30,32,30,40,38,38,34,46,40,48]$, $\max_r F(r)=48$ at $v{=}10$ & \secref{c20-twodfs} \\
0/1 pair state (1156D shape) & best $0^{*}1^{*}$ chain / crossing pair merged up and down, summed over every root & \secref{c20-updown} \\
signed prefix/suffix (960E shape) & alternating-sign path sum, $\max$-style combine not invertible; $\sum$ over ordered pairs & \secref{c20-prefixsuffix} \\
game verdict per root (1626E shape) & win/lose flag (a $\operatorname{mex}$-style value) re-evaluated for every root through the up pass & \secref{c20-counting} \\
\bottomrule
\end{tabular}
\end{center}

% ---------------- Taxonomy chart ----------------
\begin{figure}[H]
\centering
\begin{tikzpicture}[diagram, scale=0.95, transform shape,
    mnode/.style={circle, draw=#1, line width=0.9pt, fill=#1!8, minimum size=6.5mm,
                  font=\scriptsize\bfseries, text=#1!60!black}]
  \node[mnode=primaryblue] (a) at (0,0)     {$\mathit{down}[v]$};
  \node[annot, font=\scriptsize, anchor=west] at (1.15,0)  {$=$ Ch.~13's fold over the subtree; computed in the DFS that goes \emph{down} (15.1)};
  \node[mnode=successgreen] (b) at (0,-1)  {$\mathit{up}[v]$};
  \node[annot, font=\scriptsize, anchor=west] at (1.15,-1) {$=$ the parent-side piece; the edge-transfer rule moves $v$'s answer into $c$ (15.1)};
  \node[mnode=difforange]   (c) at (0,-2)  {$\mathit{dp}[v]$};
  \node[annot, font=\scriptsize, anchor=west] at (1.15,-2) {$=$ $\Phi(\mathit{down}[v],\mathit{up}[v])$, read at every node as "the answer if $v$ is root" (15.2)};
  \node[mnode=warningred]   (d) at (0,-3)  {$\text{excl}(c_i)$};
  \node[annot, font=\scriptsize, anchor=west] at (1.15,-3) {$=$ combine of all children but $c_i$, via prefix/suffix when $\Phi$ is not invertible (15.3)};
  \node[mnode=softpurple]   (e) at (0,-4)  {$g(v)$ / win};
  \node[annot, font=\scriptsize, anchor=west] at (1.15,-4) {$=$ a $\operatorname{mex}$/game value rerooted across every node (15.4, Summit)};
\end{tikzpicture}
\caption{The rerooting taxonomy -- the state splits into a down pass and an
up pass, joined by the edge-transfer rule.  Rows 1--3 are the invertible or
trivially-combinable regime; row~4 is the non-invertible regime that needs a
prefix/suffix (or exclusive-segment-tree) precompute; row~5 reroots a game
value, tying this chapter back to Chapter~14's Grundy theory.  Chapter~13's
cookbook is the "read at the root only, no up pass" special case of row~1.}
\label{fig:c20-taxonomy}
\end{figure}

% ---------------- Who owns what ----------------
\subsection*{Who Owns What (so nothing is learned twice)}

\begin{center}
\small
\renewcommand{\arraystretch}{1.25}
\begin{tabular}{@{}p{0.30\linewidth} p{0.32\linewidth} p{0.32\linewidth}@{}}
\toprule
\textbf{Idea} & \textbf{Owner} & \textbf{This chapter's share} \\
\midrule
constant-tuple child-fold, in/out pair, harvest-at-peak & Ch.~13 (\chref{c10}) & the fold \emph{is} the down pass; Ch.~15 adds the up pass and the edge-transfer rule \\
tree \emph{metrics} (depth, height, size, diameter) & Ch.~12 (\chref{c09}) & treated as already-computed inputs; $\size(v)$ drives the edge-transfer \\
budgeted $\bigO{nV}$ child-merge, game Grundy/Hackenbush & Ch.~14 (\chref{c27}) & inputs; the Summit reroots Ch.~14's game value for every root \\
answer for \emph{every root} / rerooting as a theory & \textbf{Ch.~15 (this)} & 13/14 deferred it; \thmref{c20-edge-transfer} and the two-DFS template deliver it here \\
per-pair distance \emph{queries} via LCA & Ch.~16 (\chref{c11}) & \emph{not} here -- one-shot DP per root, not $Q$ distance queries \\
small-to-large / DSU-on-tree merging & Ch.~17 (\chref{c28}) & \emph{not} here -- the up pass is a plain top-down fold \\
centroid decomposition for path-count over all roots & Ch.~24 (\chref{c18}) & when the "combine" needs a global path statistic, hand off \\
\bottomrule
\end{tabular}
\end{center}

% ---------------- Recognition guide ----------------
\begin{recognitionbox}[Pattern card -- ``the answer must be reported for every root'']
\reckey{the problem asks for a quantity $f(r)$ evaluated with the tree rooted
at \emph{each} node (a sum of subtree sizes, a total distance, a count of
valid paths/pairs through each node, a game verdict per root); or the statement
fixes no root and the sample output is an array of $n$ answers, one per node.}
\recuse{
  \item ``compute $f(r)$ for all $r$; $n$ up to $2\cdot10^{5}$'' -- re-rooting
        with the two-DFS template of \secref{c20-twodfs} (CF~1187E);
  \item the combine over children is \emph{invertible} (sum, XOR) and the
        edge-transfer is a closed form -- \secref{c20-twodfs}/\secref{c20-updown}
        subtract the child's contribution (CF~1156D);
  \item the combine is $\max$/$\min$/signed-product/$\mathit{or}$ -- \emph{not}
        invertible; use the prefix/suffix merge of \secref{c20-prefixsuffix}
        (CF~960E);
  \item a two-player game whose win/lose or Grundy value is needed at every
        root -- \secref{c20-counting} and the Summit (CF~1626E, CF~1498F).
}
\recnotuse{
  \item the answer is read at \emph{one} fixed root -- that is Chapter~13's
        cookbook (\chref{c10}); no up pass is needed;
  \item you must answer $Q$ \emph{distance/LCA queries} between named pairs --
        that is Chapter~16 (\chref{c11});
  \item you must aggregate a per-subtree multiset with heavy-child reuse --
        DSU-on-tree, Chapter~17 (\chref{c28});
  \item the quantity is a global path statistic (count paths of length $k$)
        best split at a centroid -- Chapter~24 (\chref{c18}).
}
One-line test: \emph{does the sample print one number \underline{per node}, the
value of some rooted DP as the root sweeps the whole tree?  Then it is a
rerooting problem, and this chapter is its home.}
\end{recognitionbox}

% ---------------- Constraint -> tool ----------------
\subsection*{Reading the Constraints (which rerooting shape the limits want)}

\begin{center}
\small
\renewcommand{\arraystretch}{1.25}
\begin{tabular}{@{}p{0.28\linewidth} p{0.40\linewidth} p{0.24\linewidth}@{}}
\toprule
\textbf{Constraint signature} & \textbf{What it rules in/out} & \textbf{Tool} \\
\midrule
$n\le2\cdot10^{5}$, output one value per node & $\bigO{n}$ two-DFS; a closed edge-transfer & \secref{c20-twodfs}, CF~1187E \\
combine is $+$ / $\oplus$ (invertible) & subtract the child to get "siblings", no extra precompute & \secref{c20-updown}, CF~1156D \\
combine is $\max$/$\min$/signed product ($n\le3\cdot10^{5}$) & prefix/suffix over children; $\bigO{\deg}$ per node & \secref{c20-prefixsuffix}, CF~960E \\
reroot a win/lose or Grundy value at every node & two-DFS on a $\operatorname{mex}$/flag state; bound the reachable set & \secref{c20-counting}, CF~1626E \\
game on every root with $n\le2\cdot10^{5}$, $k$ a depth cap & rerooted Hackenbush/Grundy under \thmref{c27-hackenbush} & Summit CF~1498F \\
\bottomrule
\end{tabular}
\end{center}

% ---------------- Notation ----------------
\subsection*{Notation Used Only Here}

Root the reference orientation at $1$; $c\in\mathit{ch}(v)$ ranges over the
children of $v$ in that view, and $\deg(v)$ is the number of neighbours.
Split the state: $\mathit{down}[v]$ is Ch.~13's fold over $\text{subtree}(v)$;
$\mathit{up}[v]$ is the same fold's value on the \emph{rest of the tree above
$v$} (the parent-side piece); the every-root answer is
$\mathit{dp}[v]=\Phi(\mathit{down}[v],\mathit{up}[v])$.
$\Phi$ is the operator that folds a set of child contributions into
one value -- a sum, a $\max$, a signed product, a pair-merge, a
$\operatorname{mex}$.  It is \term{invertible} when
$\Phi(S)=\Phi(S\setminus\{x\})\oplus^{-1} x$ recovers the "everything
but $c$" aggregate cheaply, and \term{non-invertible} otherwise, which forces
the prefix/suffix merge of \secref{c20-prefixsuffix}.  The \term{edge-transfer
rule} $\mathit{up}[c]=\Phi\text{-shift}(\Phi(v\text{'s self},\mathit{up}[v],\text{all siblings of }c))$
is the whole chapter; $\Phi\text{-shift}$ is whatever per-edge adjustment the
fold demands (adding $n-2\size(c)$ for CF~1187E, extending a chain by one edge
for CF~1156D, flipping a sign for CF~960E).  The seed is
$\mathit{up}[\text{root}]=\Phi\text{'s identity}$ (an empty combine),
\emph{never} a hard zero unless $\Phi$'s identity is zero.

Every section of this chapter is a chapter of the \emph{same} table on the
\emph{same} $T$: if a row here and a row there disagree, this opener is wrong
and the section is right -- fix it here, never the other way.
